\documentclass{article}
\usepackage[utf8]{inputenc}
\usepackage{amsmath}
\usepackage{geometry}
\geometry{margin=2cm}
\begin{document}

\section*{Análise da Questão 149 - ENEM 2024 - Matemática}

\subsection*{Enunciado}
Em uma empresa é comercializado um produto em embalagens em formato de cilindro circular reto, com raio medindo 3 cm e altura medindo 15 cm. Essa empresa planeja comercializar o mesmo produto em embalagens em formato de cubo, com capacidade igual a 80\% da capacidade da embalagem cilíndrica utilizada atualmente. Use 3 como valor aproximado para $\pi$. A medida da aresta da nova embalagem, em centímetro, deve ser:

\bigskip

\subsection*{Solução Técnica}

Calcula-se o volume do cilindro, que é dado por:
\[
V = \pi r^{2} h.
\]
Substituindo $r = 3$ cm, $h = 15$ cm e $\pi \approx 3$, tem-se:
\[
V = 3 \times 3^{2} \times 15 = 3 \times 9 \times 15 = 405 \text{ cm}^3.
\]
A nova embalagem terá um volume que corresponde a 80\% do volume anterior:
\[
V_{cubo} = 0,8 \times 405 = 324 \text{ cm}^3.
\]
O volume do cubo é dado por $a^{3}$, onde $a$ é a aresta. Logo,
\[
a^{3} = 324 \implies a = \sqrt[3]{324}.
\]
Fatorando 324:
\[
324 = 27 \times 12,
\]
logo,
\[
a = \sqrt[3]{27 \times 12} = \sqrt[3]{27} \times \sqrt[3]{12} = 3 \sqrt[3]{12}.
\]
A expressão $3 \sqrt[3]{12}$ corresponde à alternativa \textbf{E}.

\subsection*{Explicação Didática}

Para resolver este problema, realizamos os seguintes passos:
\begin{enumerate}
  \item Calculamos o volume da embalagem cilíndrica usando a fórmula do volume do cilindro.
  \item Calculamos 80\% deste volume para obter o volume da nova embalagem em formato de cubo.
  \item Aplicamos a fórmula do volume do cubo para encontrar a aresta a partir do volume.
  \item Calculamos a raiz cúbica do volume para determinar o comprimento da aresta.
  \item Simplificamos a expressão utilizando fatoração para facilitar a leitura e comparação com as alternativas.
\end{enumerate}

É importante destacar que a redução de volume não se traduz diretamente em redução na medida linear, pois o volume depende do cubo da medida linear (a aresta). Além disso, a raiz cúbica permite transformar o volume em comprimento da aresta.

\subsection*{Análise dos Distratores}
\begin{itemize}
  \item \textbf{Alternativa A}: Possível confusão entre raio e aresta, resultando em aresta igual a 6 cm, que não corresponde a redução de volume esperada.
  \item \textbf{Alternativa B}: Considera a altura do cilindro (18 cm, presumivelmente valor errado no enunciado) como aresta, resultando em volume maior do que o desejado.
  \item \textbf{Alternativa C}: Erro no cálculo ou simplificação inadequada da raiz cúbica, resultando em valor incorreto.
\end{itemize}

\subsection*{Perfil da Habilidade Avaliada}
Esta questão avalia habilidades de geometria espacial, cálculo de volumes, operações com radiciação e compreensão de proporcionalidade entre volumes e medidas lineares. A questão possui grau de dificuldade médio e requer cálculos precisos, além da capacidade de interpretar expressões matemáticas apresentadas.

\subsection*{Conclusão}
A alternativa correta é a \textbf{E}.

\bigskip

\textbf{Imagens das alternativas relacionadas:}

\begin{description}
  \item[Alternativa D:] $6 \sqrt[3]{6}$
  \item[Alternativa E:] $3 \sqrt[3]{12}$
\end{description}

A alternativa E, com $3 \sqrt[3]{12}$, é a que corresponde ao valor correto da aresta do cubo segundo os cálculos apresentados.

\end{document}